\emph{Preliminaries.} The blades are a basis, $e_{S}e_{S}^{\dagger} = 1$ and $\langle xy^{\dagger}\rangle_{0}$ is the sum of products of coefficients (the preliminaries of the proof of Theorem \claimref{synthesis.mode-split-of-the-clifford-algebra-and-its-klein-twist-to-the-tensor-product}); so $\langle xx^{\dagger}\rangle_{0} \ge 0$, $\langle ab\rangle_{0} = \langle ba\rangle_{0}$, and a self-adjoint blade has square $1$. For commuting pairs: $(A_{i}B_{j})^{\dagger} = B_{j}A_{i} = A_{i}B_{j}$, so $C^{\dagger} = C$; with $S = B_{1} + B_{2}$, $D = B_{1} - B_{2}$, $C = A_{1}S + A_{2}D$ and $C^{2} = S^{2} + D^{2} + A_{1}A_{2}SD + A_{2}A_{1}DS = 4 - [A_{1},A_{2}][B_{1},B_{2}]$, since $SD = -[B_{1},B_{2}] = -DS$. For a state, $0 \le \varphi((C - t)(C - t)^{\dagger})$ at $t = \varphi(C)$ gives $\varphi(C)^{2} \le \varphi(C^{2})$.

\emph{(i).} Clause (i) of the stored theorem gives $xy = (-1)^{|x||y|}yx$ for homogeneous $x \in \Cl(A)$, $y \in \Cl(B)$, so even elements commute. Clause (ii) gives $y\Gamma_{A}y'\Gamma_{A} = yy'$ for generators of $\Cl(B)$ and that the even part of $B'$ is the even part of $\Cl(B)$. An element of $\Cl(A)$ or of $\Cl(B)$ commutes with $\Gamma$ exactly when it is even, since $\Gamma$ anticommutes with every generator.

\emph{(ii).} The blades of the algebra of mode $k$ are $1$, $e_{k}$, $f_{k}$, $p_{k} = f_{k}e_{k}$; $f_{k}^{\dagger} = -f_{k}$ and $p_{k}^{\dagger} = e_{k}^{\dagger}f_{k}^{\dagger} = -e_{k}f_{k} = p_{k}$, so the self-adjoint elements are $s + tp_{k} + ue_{k}$. Since $p_{k}$ and $e_{k}$ anticommute and have square $1$, $(s + tp_{k} + ue_{k})^{2} = s^{2} + t^{2} + u^{2} + 2s(tp_{k} + ue_{k})$, which is $1$ exactly when $t = u = 0$, $s = \pm 1$, or $s = 0$, $t^{2} + u^{2} = 1$. The element is even exactly when $u = 0$. So the even observables are $\pm 1$, $\pm p_{k}$; they commute, $[A_{1},A_{2}] = 0$, $C^{2} = 4$ and $\varphi(C)^{2} \le 4$. Nothing was assumed about the other party beyond $A_{i}B_{j} = B_{j}A_{i}$.

\emph{(iii).} $e_{1}f_{2}$ is self-adjoint ($f_{2}^{\dagger}e_{1}^{\dagger} = -f_{2}e_{1} = e_{1}f_{2}$) with square $-e_{1}^{2}f_{2}^{2} = 1$, and it anticommutes with $p_{1} = f_{1}e_{1}$, which contains $e_{1}$ and not $f_{2}$; likewise for $B$. All four are products of two generators, hence even, and commute across the parties by (i). Then $[A_{1},A_{2}] = 2p_{1}e_{1}f_{2} = 2f_{1}f_{2}$, $[B_{1},B_{2}] = 2f_{3}f_{4}$ and $C^{2} = 4 - 4f_{1}f_{2}f_{3}f_{4} = 4 + 4K$. $K$ is a product of the commuting anti-self-adjoint elements $f_{1}f_{2}$, $f_{3}f_{4}$ of square $-1$, so $K^{\dagger} = K$, $K^{2} = 1$; it is even, and it commutes with each $A_{i}B_{j}$ because $A_{i}$ anticommutes with $A_{1}A_{2}$ and $B_{j}$ with $B_{1}B_{2}$. As the $A$'s commute with the $B$'s, $KA_{1}B_{1} = -(A_{1}A_{2}A_{1})(B_{1}B_{2}B_{1}) = -A_{2}B_{2}$, $KA_{2}B_{2} = -A_{1}B_{1}$, $KA_{1}B_{2} = A_{2}B_{1}$, $KA_{2}B_{1} = A_{1}B_{2}$, so $KC = C$ and $C^{3} = 8C$. The four terms of $C$ and $K$ are blades other than $1$, so $\langle C\rangle_{0} = \langle KC\rangle_{0} = \langle K\rangle_{0} = 0$, whence $\varphi_{r}(1) = 1$, $\varphi_{r}(C) = r\langle(1 + K)(4 + 4K)\rangle_{0} = 8r$ and $\varphi_{r}(C^{2}) = 8 + 8r\langle(1 + K)C\rangle_{0} = 8$. Positivity: with $P = \tfrac12(1 + K)$ and $y = Px$, $\tfrac12\varphi_{r}(xx^{\dagger}) = \langle y^{\dagger}y\rangle_{0} + r\langle y^{\dagger}Cy\rangle_{0}$, and by the Cauchy-Schwarz inequality for the positive form $\langle a^{\dagger}b\rangle_{0}$, $\langle y^{\dagger}Cy\rangle_{0}^{2} \le \langle y^{\dagger}y\rangle_{0}\langle y^{\dagger}C^{2}y\rangle_{0} = 8\langle y^{\dagger}y\rangle_{0}^{2}$ (as $C^{2}P = 8P$); so $8r^{2} \le 1$ gives $\varphi_{r}(xx^{\dagger}) \ge 0$. The element $(1 + K)(1 + rC)$ is even, so $\varphi_{r}$ is even. The rational $r$ with $8r^{2} \le 1$ have supremum $\sqrt{2}/4$, not rational.

\emph{(iv).} By clause (ii) of the stored theorem $e_{2}p_{1}$ has square $1$ and commutes with $\Cl(A)$; it is self-adjoint since $e_{2}$ and $p_{1}$ are and commute; it anticommutes with $p_{2}$ as $e_{2}$ does; $p_{2}$ commutes with $\Cl(A)$. So the four elements are a CHSH configuration. The state is the pair state of clause (v) of the stored theorem with $c = s$, $r = 0$, there shown to be pure and even. Expanding the null vectors, $w_{1}w_{2} + v_{2}v_{1} = -\tfrac12(e_{1}f_{2} + f_{1}e_{2})$, so $\rho = \tfrac14(1 + p_{1}p_{2} - e_{1}f_{2} - f_{1}e_{2})$, a sum of four blades of square $1$. The value $4\langle\rho x\rangle_{0}$ of a blade $x$ of square $1$ is its coefficient in $4\rho$: it is $1$ for $p_{1}p_{2}$ and for $e_{1}\,e_{2}p_{1} = e_{2}f_{1} = -f_{1}e_{2}$, and $0$ for the odd blades $p_{1}\,e_{2}p_{1} = e_{2}$ and $e_{1}p_{2}$. With $B_{\pm} = cp_{2} \pm s\,e_{2}p_{1}$: $\varphi(C) = \varphi(p_{1}(B_{+} + B_{-})) + \varphi(e_{1}(B_{+} - B_{-})) = 2c\,\varphi(p_{1}p_{2}) + 2s\,\varphi(e_{1}e_{2}p_{1}) = 2(c + s)$, and $(c + s)^{2} = 1 + 2cs > 1$ for $c, s > 0$. $\Gamma$ anticommutes with every generator, so with $e_{1}$ and with $e_{2}p_{1}$ (three generators), and the parity average $E(x) = \tfrac12(x + \Gamma x\Gamma)$, which is the average of Theorem \claimref{singularity.descent-parity-average} with the total parity $\Gamma$ in place of its single-mode $\sigma$ (the form used in the stored theorem on the mode split), sends them to $0$; $E(p_{k}) = p_{k}$; the averaged CHSH element is $2c\,p_{1}p_{2}$, of value $2c$. The last statement is (ii) applied to either party: in the twisted split the observables of $B'$ are $\pm 1$ and $tp_{2} + u\,e_{2}p_{1}$ by the isomorphism $B' \cong \Cl(1,1)$, even exactly when $u = 0$.

\emph{(v).} $e_{1}e_{2} = -e_{2}e_{1}$ and $(e_{1}e_{2})^{\dagger} = e_{2}e_{1} = -e_{1}e_{2}$. If $\varphi(x^{\dagger}) = \varphi(x)$ for all $x$ then $\varphi(e_{1}e_{2}) = -\varphi(e_{1}e_{2}) = 0$; for $\varphi = \langle\rho\,\cdot\,\rangle_{0}$ with $\rho^{\dagger} = \rho$ this holds because $\langle\rho x^{\dagger}\rangle_{0} = \langle(x\rho)^{\dagger}\rangle_{0} = \langle x\rho\rangle_{0} = \langle\rho x\rangle_{0}$. (Over real coefficients positivity alone does not give $\varphi(x^{\dagger}) = \varphi(x)$, which is why it is a hypothesis of the clause.) Put $X = e_{1}p_{2}$, $Y = p_{1}e_{2}$, $Z = p_{1}p_{2}$: each is a product of two commuting self-adjoint elements of square $1$, so is self-adjoint of square $1$, and $(e_{1}e_{2})$ is the only term of the untwisted element that is not self-adjoint, so its self-adjoint part is $H = X + Y - Z$. $XZ = e_{1}p_{1} = -p_{1}e_{1} = -ZX$, likewise $YZ = -ZY$, and $XY = e_{1}p_{1}\,p_{2}e_{2}$, $YX = p_{1}e_{2}e_{1}p_{2} = -p_{1}e_{1}\,e_{2}p_{2} = -XY$. So $H^{2} = 3$ and $\varphi(H)^{2} \le \varphi(H^{2}) = 3$ for every state.
